\documentclass[11pt,a4paper]{article}
\usepackage{fontspec,amsmath,amssymb,amsthm,microtype}
\setmainfont{lmroman10-regular.otf}[BoldFont=lmroman10-bold.otf,ItalicFont=lmroman10-italic.otf,BoldItalicFont=lmroman10-bolditalic.otf]
\usepackage[margin=25mm]{geometry}
\usepackage[unicode,hidelinks]{hyperref}
\hypersetup{pdfauthor={LangevinLean Project}}
\setlength{\emergencystretch}{3em}
\setlength{\parskip}{0.4em}
\setlength{\parindent}{0pt}
\allowdisplaybreaks[2]
\raggedbottom
\title{\texorpdfstring{Contraction of shared-\allowbreak{}input kinetic integral solutions}{Contraction of shared-input kinetic integral solutions}}
\author{LangevinLean Project}
\date{}
\begin{document}\maketitle
Draft mathematical exposition; these notes may contain mistakes.


This joins the Picard integral construction to the transformed-energy argument. Two solutions with the same continuous forcing have a noise-free difference equation as a consequence of their integral identities. The energy estimate and the paper's position decay then follow without assuming either a contraction result or a differentiable Brownian path.



The abstract energy inequalities are applied to solutions already constructed from their inputs. This closes the gap between a formal difference equation and a contraction statement for kinetic trajectories.


\section{\texorpdfstring{Derive the phase difference equation and contract its energy}{Derive the phase difference equation and contract its energy}}


Apply the shared-input derivative theorem to the concrete kinetic drift \(G(v,x)=(\gamma v+\nabla f(x),-v)\). Project its right derivative to velocity and position using the two continuous linear projection maps. The velocity derivative becomes \(-\gamma(V-V\prime)-[\nabla f(X)-\nabla f(X\prime)]\); the position derivative is \(V-V\prime\). Scalar multiplication and additive rearrangement justify both identities. The right-derivative energy theorem, which is valid at the initial endpoint, yields

\[\mathcal E_t\le e^{-2mt/\gamma}\mathcal E_0,\quad\mathcal E_t=\|V_t-V\prime_t+\gamma(X_t-X\prime_t)\|^2+\|V_t-V\prime_t\|^2.\]

In The proof the required Lipschitz constant for \(G\) is derived from the gradient bound; it is not an extra Hessian or solution assumption. The argument then supplies the original potential's strong convexity and smoothness hypotheses to the proved energy inequality.



The contraction theorem is now applied to integral solutions. Its differential hypotheses follow by cancellation and the fundamental theorem, so they are conclusions of the construction rather than extra assumptions about an arbitrary pair of random paths.


\section{\texorpdfstring{Equal initial velocities give the sharp position prefactor}{Equal initial velocities give the sharp position prefactor}}


If the initial velocities agree, \(\mathcal E_0=\gamma^2\|X_0-X\prime_0\|^2\). The projection estimate \(\gamma^2\|X_t-X\prime_t\|^2\le2\mathcal E_t\) cancels the factor \(\gamma^2\), giving the squared position bound. Positivity of \(\gamma\), the norm, the exponential and the square root permits taking square roots. The argument explicitly proves \(e^{-mt/\gamma\,\cdot2}=(e^{-mt/\gamma})^2\), rewrites \((\sqrt2)^2=2\), and applies the nonnegative square comparison theorem. The resulting prefactor is \(\sqrt2\), consistent with the explicit general prefactor in the supplied source.



The initial velocity difference is zero because the coupling shares one initial velocity. Thus its initial phase cost is exactly a constant times the position discrepancy. This special initialization is why the final theorem can be stated using position laws alone.


\section{\texorpdfstring{Contract the full transformed norm}{Contract the full transformed norm}}


The integral-flow energy theorem gives \(\mathcal E_t\le e^{-2mt/\gamma}\mathcal E_0\). Replacing energy by the square of the transformed norm and using \((e^{-mt/\gamma})^2=e^{-2mt/\gamma}\) gives \(\|Y_t-Z_t\|_T\le e^{-mt/\gamma}\|Y_0-Z_0\|_T\). One proves the exponential square identity by the exponential addition formula, notes that both norms and the exponential are nonnegative, and invokes the energy-to-norm lemma. This stronger phase estimate controls velocity differences and accumulates discretization errors without a separate projection loss at each step.



Unlike the equal-velocity projection, contraction in the full transformed norm permits any initial phase difference. This stronger form is needed at later numerical steps, where an accumulated velocity error is generally present as well as a position error.


\section{\texorpdfstring{Apply the full shifted theorem to integral flows}{Apply the full shifted theorem to integral flows}}


Take two kinetic integral solutions with the same additive input. Their noise-cancelled phase difference has a right derivative by the right derivative of the subtracted integral equations. Compose this derivative with the continuous linear velocity and position projections; vector-space algebra normalizes their drift formulas. The centered-sector theorem then supplies the full arbitrary-shift energy exponent from the original potential assumptions. Equal initial velocities identify the exact initial energy, and the position-square projection supplies the final prefactor. Squaring the desired nonnegative quantities, using the exponential-addition identity, and dividing by \(\gamma-2a>0\) proves

\[\|X_t-Y_t\|\le\frac{\sqrt{2((\gamma-a)^2+a^2)}}{\gamma-2a}\exp\!\left(\frac{\max\{a^2-m,M-(\gamma-a)^2\}}{\gamma-2a}t\right)\|X_0-Y_0\|.\]

No continuous-Hessian assumption or difference-ODE assumption remains; both the ODE and the coefficient are derived from the concrete integral paths.



The shift controls the relative weighting of the two phase components. The same integral-solution argument works for every admissible shift, so optimization changes only the scalar prefactor and exponent. The stochastic construction need not be repeated for each friction regime.


\end{document}
